\section{Transparency: can I see why it did that?}
\label{sec:transparency}

A clinician cannot endorse a black box. The third question is whether the system's
reasoning is inspectable and whether a complete, tamper-evident record lets the
clinician reconstruct events after the fact. Transparency is the difference between
a tool a clinician can defend and one a clinician can only hope was right. The
relevant regulatory direction is explicit on this point: the FDA's action plan for
artificial-intelligence and machine-learning software as a medical device places
transparency to users and traceability of decisions at the center of trustworthy
deployment \cite{fda2021samd}.

The mechanism makes transparency operational rather than rhetorical. Because every
action passes through verification before generation and a ten-gate VVUQ
examination resolving to ACCEPT, ESCALATE, or BLOCK, each of those steps produces a
durable artifact. Together those artifacts form an audit trail: an append-only,
hash-chained record in which each entry is cryptographically linked to the one
before it, so that any later alteration breaks the chain and is detectable. The
clinician is not asked to trust that the record is complete; the clinician can
verify it. The author's records-and-server lineage describes the national MCP
servers that hold these records under a common schema
\cite{Kawchak2026NationalMCPServersPhysical}, and the site documentation package
specifies exactly which elements are captured at each step and how they are bound
together \cite{Kawchak2026PhysicalAIOncologyClinical}.

\autoref{tab:transparency} lists what the record contains and who reads each part,
because transparency means different things to different readers. The clinician
needs the model's reasoning and the gate results to make and defend a decision; an
auditor needs the inputs and the hash chain to confirm that nothing was altered
after the fact; a tumor board needs the countersignature that ties a human to the
action. Figure 5 shows the same record being written as the action proceeds and
then read back by an auditor who reconstructs the event from the record alone.

This is what lets a clinician answer the two questions that decide whether a tool
survives contact with real practice. Can I defend this decision at tumor board?
Yes, because the model's reasoning and the gate results are on the record, not in
anyone's memory. If an auditor asked me tomorrow, could I show exactly what happened
and why? Yes, because the inputs, the reasoning, the gate outcomes, and my
countersignature are bound into a hash chain that an outside party can replay
end to end. A record that can be reconstructed by someone who was not in the room is
the test the FDA direction implies and the one this framework is built to pass.

\begin{cfwfig}
\node[cftwo] (a) at (0,0) {Action}; \node[cfact] (b) at (3.6,0) {Hash-chained record};
\node[cfhope] (c) at (7.6,0) {Auditor reconstructs};
\draw[cfedge] (a)--(b); \draw[cfedge] (b)--(c);
\draw[cfedged] (b) to[out=120,in=60] (a.north);
\end{cfwfig}
\vspace{-0.2cm}
\cfwcaption{Figure 5. The audit trail (Mermaid 08 reproduced as colored TikZ): each
action writes a hash-chained record from which an auditor can later reconstruct
exactly what happened and why.}

\needspace{9\baselineskip}\begin{table}[H]
\footnotesize
\begin{tabularx}{\textwidth}{@{}L{3.2cm} Y L{3.2cm}@{}}
\toprule
\textbf{Record element} & \textbf{Contents} & \textbf{Who reads it}\\
\midrule
Inputs & The patient data, parameters, and context presented to the system & Auditor, sponsor\\
Model reasoning & The proposed action and the explanation Claude Code generated for it & Clinician, tumor board\\
Gate results & The pass, ESCALATE, or BLOCK outcome of each of the ten VVUQ gates & Clinician, auditor\\
Clinician countersignature & The named human who accepted, escalated, or overrode the action & Tumor board, IRB\\
Hash chain & The cryptographic link binding every entry into a tamper-evident sequence & Auditor, regulator\\
\bottomrule
\end{tabularx}
\caption{The audit-trail record: what each element contains and which reader relies
on it.}
\label{tab:transparency}
\end{table}

Transparency is operational, not rhetorical: the test is whether an auditor can
reconstruct exactly what happened, and why, from the record alone. The hash-chained
audit trail is that record, and it is the concrete provision that turns a black box
into a tool a clinician can defend at tumor board.
